\documentclass[../Main.tex]{subfiles}
\begin{document}
  \chapter{2020-2021学年微积分（一）（上）期中考试}

  \section{单项选择题(每小题 3 分，共 18 分)}
  \begin{enumerate}
    \item 设数列$\{x_{n}\}$收敛，则【\quad】.

      \begin{tasks}[label=\textbf{\Alph*.}, label-width=1.5em, column-sep=1em](1)
        \task 当$\lim_{n\to\infty}\sin x_{n}=0$时，$\lim_{n\to\infty}x_{n}=0$
        \task 当$\lim_{n\to\infty}\left(x_{n}+\sqrt{|x_{n}|}\right)=0$时，$\lim_{n\to\infty}x_{n}=0$
        \task 当$\lim_{n\to\infty}\left(x_{n}+x_{n}^{2}\right)=0$时，$\lim_{n\to\infty}x_{n}=0$
        \task 当$\lim_{n\to\infty}\left(x_{n}+\sin x_{n}\right)=0$时，$\lim_{n\to\infty}x_{n}=0$
      \end{tasks}

    \item 极限$\lim_{x\to\infty}\left[\frac{x^{2}}{(x-a)(x+b)}\right]^{x}$=【\quad】.

      \begin{tasks}[label=\textbf{\Alph*.}, label-width=1.5em, column-sep=1em](2)
        \task $1$
        \task $\mathrm{e}$
        \task $\mathrm{e}^{a-b}$
        \task $\mathrm{e}^{b-a}$
      \end{tasks}

    \item 设$F(x)=
      \begin{cases}
        \frac{f(x)}{x}, & x\neq0, \\
        f'(0),          & x=0
      \end{cases}$，其中$f'(0)=1$，$f(0)=0$，则$x=0$是$F(x)$的【\quad】.

      \begin{tasks}[label=\textbf{\Alph*.}, label-width=1.5em, column-sep=1em](2)
        \task 连续点
        \task 可去间断点
        \task 无穷间断点
        \task 跳跃间断点
      \end{tasks}

    \item 函数$f(x)$在$x=0$处可导，且$f(0)=0$，则$\lim_{x\to0}\frac{x^{2}f(x)-2f(x^{3})}{x^{3}}$=【\quad】.

      \begin{tasks}[label=\textbf{\Alph*.}, label-width=1.5em, column-sep=1em](2)
        \task $-2f'(0)$
        \task $-f'(0)$
        \task $f'(0)$
        \task $0$
      \end{tasks}

    \item 设$f(x)$具有一阶连续导数，$F(x)=f(x)\left(1+|\sin x|\right)$，则$f(0)=0$是$F
      (x)$在$x=0$处可导的【\quad】.

      \begin{tasks}[label=\textbf{\Alph*.}, label-width=1.5em, column-sep=1em](2)
        \task 充分必要条件
        \task 充分条件但非必要条件
        \task 必要条件但非充分条件
        \task 既非充分又非必要条件
      \end{tasks}

    \item 设$f(x)$有二阶连续导数，且$f'(0)=0$，$\lim_{x\to0}\frac{f''(x)}{|x|}=1$，则【\quad】.

      \begin{tasks}[label=\textbf{\Alph*.}, label-width=1.5em, column-sep=1em](2)
        \task $f(0)$是$f(x)$的极大值
        \task $f(0)$是$f(x)$的极小值
        \task $\left(0,f(0)\right)$是曲线$y=f(x)$的拐点
        \task $f'(0)$是$f'(x)$的极值
      \end{tasks}
  \end{enumerate}

  \section{填空题(每小题 4 分，共 16 分)}
  \begin{enumerate}
    \setcounter{enumi}{6}
    \item 极限$\lim_{x\to0}x\left[\frac{1}{x}\right]=\underline{\hspace{3cm}}$.

    \item 设$y=\frac{\cos x}{x^{2}}$，则$\frac{\mathrm{d}y}{\mathrm{d}(\cos
      x)}=\underline{\hspace{3cm}}$.

    \item $\lim_{x\to0}\left(\frac{1}{\mathrm{e}^{x}-1}-\frac{1}{\ln(1+x)}\right
      )=\underline{\hspace{3cm}}$.

    \item 曲线$y=x\left(1+\arcsin\frac{2}{x}\right)$的斜渐近线方程为\underline{\hspace{3cm}}.
  \end{enumerate}

  \section{基本计算题(每小题 7 分，共 42 分)}
  \begin{enumerate}
    \setcounter{enumi}{10}
    \item 设$f(x)=x+a\ln(1+x)+bx\sin x$，$g(x)=cx^{3}$，在$x\to0$时，若$f(x)$与$g
      (x)$是等价无穷小，求常数$a,b,c$的值.
      \vspace{10em}

    \item 设函数$f(x)=
      \begin{cases}
        x\arctan\frac{1}{x^{2}}, & x\neq0, \\
        0,                       & x=0
      \end{cases}$，讨论$f'(x)$在$x=0$处的连续性.
      \vspace{10em}

    \item 设$\begin{cases}
        x=\sin t,        \\
        y=t\sin t+\cos t
      \end{cases}$，$t\in\left(0,\frac{\pi}{2}\right)$，求$\left.\frac{\mathrm{d}^{2}
      y}{\mathrm{d}x^{2}}\right|_{t=\frac{\pi}{4}}$.
      \vspace{10em}

    \item 设$y=\frac{x^{2}}{2-x}$，计算$y^{(50)}(0)$.
      \vspace{10em}

    \item 设$x=\varphi(y)$是函数$y=x\ln x$的反函数，计算$\varphi(y)$在$x=\mathrm{e}$处的二阶导数$\frac{\mathrm{d}^{2}
      x}{\mathrm{d}y^{2}}$.
      \vspace{12em}

    \item 求$f(x)=x^{3}-3x^{2}-9x+5$在区间$[-2,2]$上的最大值与最小值.
      \vspace{12em}
  \end{enumerate}

  \section{应用题(每小题 7 分，共 14 分)}
  \begin{enumerate}
    \setcounter{enumi}{16}
    \item 如果以每秒$50\mathrm{cm}^{3}$的速度匀速给一个气球充气，假设气球内气压保持常值且形状始终为球形，问当气球的半径为$5
      \mathrm{cm}$时，气球半径增加的速率是多少？
      \vspace{15em}

    \item 设在$\left(-\infty,+\infty\right)$内$f''(x)>0$，$f(0)<0$，讨论$F(x
      )=\frac{f(x)}{x}(x\neq0)$的单调性.
      \vspace{15em}
  \end{enumerate}

  \section{综合题(每小题 5 分，共 10 分)}
  \begin{enumerate}
    \setcounter{enumi}{18}
    \item 设$x_{n}=\sqrt{a+\sqrt{a+\cdots+\sqrt{a}}}$（$n$个根式，$a>0$），证明$\{
      x_{n}\}$收敛，并求出极限值.
      \vspace{18em}

    \item 设函数$f(x)$在区间$[0,2]$上具有连续导数，$f(0)=f(2)=0$，$M=\max_{x\in(0,2)}
      \{|f(x)|\}$. 证明：$\exists\,\xi\in(0,2)$，使得$|f'(\xi)|\geq M$.
  \end{enumerate}

  \chapter{2020-2021学年微积分（一）（上）期中考试参考答案}
  \section{单项选择题 (每小题 3 分，共 18 分)}
  \begin{enumerate}
    \item \textbf{Solution}. D.

      设$\lim_{n\to\infty}x_{n}=A$，分别解方程$\sin A=0$，$A+\sqrt{|A|}=0$，$A+A^{2}
      =0$，$A+\sin A=0$，

      只有$A+\sin A=0$具有唯一解$A=0$.

    \item \textbf{Solution}. C.
      \[
        \begin{aligned}
          \lim_{x\to\infty}\left[\frac{x^2}{(x-a)(x+b)}\right]^{x}= & \lim_{x\to\infty}\left[1+\frac{(a-b)x+ab}{(x-a)(x+b)}\right]^{x}                                                                    \\
          =                                                         & \lim_{x\to\infty}\left[1+\frac{(a-b)x+ab}{(x-a)(x+b)}\right]^{\frac{(x-a)(x+b)}{(a-b)x+ab}\cdot\frac{(a-b)x+ab}{(x-a)(x+b)}\cdot x} \\
          =                                                         & \mathrm{e}^{\frac{(a-b)x^2+abx}{x^2+(b-a)x-ab}}                                                                                     \\
          =                                                         & \mathrm{e}^{a-b}.
        \end{aligned}
      \]

    \item \textbf{Solution}. B.

      因
      \[
        \lim_{x\to0}F(x)=\lim_{x\to0}\frac{f(x)}{x}=\lim_{x\to0}\frac{f(x)-f(0)}{x-0}
        =f'(0)=1,
      \]

      且$F(0)=0$，所以$x=0$是$F(x)$的可去间断点.

    \item \textbf{Solution}. B.

      \[
        \begin{aligned}
          \lim_{x\to0}\frac{x^2f(x)-2f(x^3)}{x^3}= & \lim_{x\to0}\left[\frac{f(x)-f(0)}{x}-2\frac{f(x^3)-f(0)}{x^3}\right] \\
          =                                        & f'(0)-2f'(0)=-f'(0).
        \end{aligned}
      \]

    \item \textbf{Solution}. A.

      若$f(0)=0$，$F(0)=0$，则
      \[
        \begin{aligned}
          \lim_{x\to0}\frac{F(x)-F(0)}{x-0}= & \lim_{x\to0}\frac{f(x)(1+|\sin x|)}{x}                    \\
          =                                  & \lim_{x\to0}f'(0)\cdot\lim_{x\to0}\left(1+|\sin x|\right) \\
          =                                  & f'(0),
        \end{aligned}
      \]

      即$F(x)$在$x=0$处可导.

      若$F(x)$在$x=0$处可导，$\lim_{x\to0}\frac{F(x)-F(0)}{x}$存在.又
      \[
        \begin{aligned}
          \lim_{x\to0}\frac{F(x)-F(0)}{x}= & \lim_{x\to0}\frac{f(x)(1+|\sin x|)-f(0)}{x}                                         \\
          =                                & \lim_{x\to0}\frac{f(x)-f(0)}{x}+\lim_{x\to0}f(x)\cdot\lim_{x\to0}\frac{|\sin x|}{x} \\
          =                                & f'(0)+f(0)\cdot\lim_{x\to0}\frac{|\sin x|}{x}.
        \end{aligned}
      \]

      由于$\lim_{x\to0}\frac{|\sin x|}{x}$不存在，所以必须$f(0)=0$.

      因此$f(0)=0$是$F(x)$在$x=0$处可导的充分必要条件.

    \item \textbf{Solution}. B.

      因$\lim_{x\to0}\frac{f''(x)}{|x|}=1$，由极限的保号性知存在$\delta>0$，当$x\in
      (-\delta,\delta)$时，$f''(x)\geq0$恒成立.

      利用Taylor展开，
      \[
        f(x)=f(0)+f'(0)x+\frac{f''(\xi)}{2}x^{2}=f(0)+\frac{f''(\xi)}{2}x^{2} \geq
        f(0), \quad -\delta<x<\delta,
      \]

      其中$\xi$介于$0$与$x$之间.故$f(0)$是$f(x)$的极小值点.
  \end{enumerate}

  \section{填空题(每小题 4 分，共 16 分)}
  \begin{enumerate}
    \setcounter{enumi}{6}
    \item \textbf{Solution}. $0$.

      利用$\frac{1}{x}\leq\left[\frac{1}{x}\right]\leq\frac{1}{x}+1$，且$\lim_{x\to0}
      x\cdot\frac{1}{x}=1=\lim_{x\to0}x\left(\frac{1}{x}+1\right)$，

      由夹逼定理可知$\lim_{x\to0}x\left[\frac{1}{x}\right]=1$.

    \item \textbf{Solution}. $\frac{x\sin x+2\cos x}{x^{3}\sin x}$.

      由链式法则知
      \[
        \begin{aligned}
          \frac{\mathrm{d}y}{\mathrm{d}(\cos x)} = & \frac{\frac{\mathrm{d}y}{\mathrm{d}x}}{\frac{\mathrm{d}(\cos x)}{\mathrm{d}x}} \\
          =                                        & \frac{\frac{-x^{2}\sin x - 2x\cos x}{(x^{2})^{2}}}{-\sin x}                    \\
          =                                        & \frac{x\sin x+2\cos x}{x^{3}\sin x}.
        \end{aligned}
      \]

    \item \textbf{Solution}. $-1$.
      \[
        \begin{aligned}
          \lim_{x\to0}\left(\frac{1}{\mathrm{e}^x-1}-\frac{1}{\ln(1+x)}\right) = & \lim_{x\to0}\frac{\ln(1+x)-(\mathrm{e}^x-1)}{(\mathrm{e}^x-1)\ln(1+x)}                          \\
          =                                                                      & \lim_{x\to0}\frac{\left(x-\frac{x^2}{2}+o(x^2)\right)-\left(x+\frac{x^2}{2}+o(x^2)\right)}{x^2} \\
          =                                                                      & \lim_{x\to0}\frac{-x^2+o(x^2)}{x^2}                                                             \\
          =                                                                      & -1.
        \end{aligned}
      \]

    \item \textbf{Solution}. $y=x+2$.

      $k=\lim_{x\to\infty}\frac{x\left(1+\arcsin\frac{2}{x}\right)}{x}=\lim_{x\to\infty}
      \left(1+\arcsin\frac{2}{x}\right)=1$，

      $b=\lim_{x\to\infty}\left[x\left(1+\arcsin\frac{2}{x}\right)-x\right]=\lim_{x\to\infty}
      x\arcsin\frac{2}{x}=2$.

      故斜渐近线方程为$y=x+2$.
  \end{enumerate}

  \section{基本计算题(每小题 6 分，共 30 分)}
  \begin{enumerate}
    \setcounter{enumi}{10}
    \item \textbf{Solution}. 利用Taylor公式，
      \[
        \begin{aligned}
          f(x)= & x+a\left(x-\frac{x^2}{2}+\frac{x^3}{3}+o(x^{3})\right)+bx\left(x-\frac{x^3}{6}+o(x^{3})\right) \\
          =     & (1+a)x+\left(b-\frac{a}{2}\right)x^{2}+\frac{a}{3}x^{3}+o(x^{3}).
        \end{aligned}
      \]

      因$f(x)$与$g(x)$在$x\to0$时是等价无穷小，故$1+a=0$，$b-\frac{a}{2}=0$，$\frac{a}{3}
      =c$.

      解得$a=-1$，$b=-\frac{1}{2}$，$c=-\frac{1}{3}$.

    \item \textbf{Solution}. 当$x\neq0$时，$f'(x)=\arctan\frac{1}{x^{2}}-\frac{2x^{2}}{x^{4}+1}$.

      当$x=0$时，$f'(0)=\lim_{x\to0}\frac{f(x)-f(0)}{x-0}=\lim_{x\to0}\arctan\frac{1}{x^{2}}
      =\frac{\pi}{2}$.

      因$\lim_{x\to0}\left(\arctan\frac{1}{x^{2}}-\frac{2x^{2}}{x^{4}+1}\right)=\lim
      _{x\to0}\arctan\frac{1}{x^{2}}=\frac{\pi}{2}=f'(0)$，所以$f'(x)$在$x=0$处连续.

    \item \textbf{Solution}. $\frac{\mathrm{d}y}{\mathrm{d}x}=\frac{\frac{\mathrm{d}y}{\mathrm{d}t}}{\frac{\mathrm{d}x}{\mathrm{d}t}}
      =\frac{t\cos t}{\cos t}=t$.

      $\frac{\mathrm{d}^{2} y}{\mathrm{d}x^{2}}=\frac{\mathrm{d}}{\mathrm{d}x}\left
      (\frac{\mathrm{d}y}{\mathrm{d}x}\right)=\frac{\frac{\mathrm{d}}{\mathrm{d}t}\left(\frac{\mathrm{d}y}{\mathrm{d}x}\right)}{\frac{\mathrm{d}x}{\mathrm{d}t}}
      =\frac{1}{\cos t}$.

      当$t=\frac{\pi}{4}$时，$\left.\frac{\mathrm{d}^{2} y}{\mathrm{d}x^{2}}\right
      |_{t=\frac{\pi}{4}}=\sqrt{2}$.

    \item \textbf{Solution}. $y=\frac{x^{2}}{2-x}=\frac{(2-x)^{2}-4(2-x)+4}{2-x}
      =\frac{4}{2-x}-2-x$.

      利用$\left(\frac{1}{a-x}\right)^{(n)}=\frac{n!}{(a-x)^{n+1}}$，得
      \[
        y^{(50)}(0)=\left.\left(\frac{4}{2-x}\right)^{(50)}\right|_{x=0}-0=\frac{4\cdot50!}{2^{51}}
        =\frac{50!}{2^{49}}.
      \]

    \item \textbf{Solution}. $\frac{\mathrm{d}x}{\mathrm{d}y}=\frac{1}{\ln
      x+1}$，$\frac{\mathrm{d}^{2} x}{\mathrm{d}y^{2}}=\frac{\mathrm{d}}{\mathrm{d}y}
      \left(\frac{1}{\ln x+1}\right)=\frac{\mathrm{d}}{\mathrm{d}x}\left(\frac{1}{\ln
      x+1}\right)\cdot\frac{\mathrm{d}x}{\mathrm{d}y}=-\frac{1}{x(\ln x+1)^{3}}$.

      当$x=\mathrm{e}$时，$\frac{\mathrm{d}^{2} x}{\mathrm{d}y^{2}}=-\frac{1}{8\mathrm{e}}$.

    \item \textbf{Solution}. $f'(x)=3x^{2}-6x-9=3(x+1)(x-3)$，故$f(x)$在$[-2
      ,2]$内的极值点为$x=-1$.

      计算$f(-2)=3$，$f(-1)=10$，$f(2)=-17$.

      故$f(x)$在$[-2,2]$上的最大值为$10$，最小值为$-17$.
  \end{enumerate}

  \section{应用题(每小题 7 分，共 14 分)}
  \begin{enumerate}
    \setcounter{enumi}{16}
    \item \textbf{Solution}. 设气球的半径为$r$，体积为$V$，则$V=\frac{4}{3}\pi r
      ^{3}$.

      由题意得$\frac{\mathrm{d}V}{\mathrm{d}t}=50$，故
      \[
        \frac{\mathrm{d}r}{\mathrm{d}t}=\frac{\mathrm{d}V}{\mathrm{d}t}\cdot\frac{\mathrm{d}r}{\mathrm{d}V}
        =\frac{50}{4\pi r^{2}}.
      \]

      当$r=5\mathrm{cm}$时，$\frac{\mathrm{d}r}{\mathrm{d}t}=\frac{50}{100\pi}=\frac{1}{2\pi}
      \mathrm{cm/s}$.

    \item \textbf{Solution}. $F'(x)=\frac{xf'(x)-f(x)}{x^{2}}$，令$G(x)=xf'(x)-
      f(x)$，则$G'(x)=xf''(x)$.

      由$f''(x)>0$知$G'(x)=0$具有唯一解$x=0$，且$G(x)$在$(-\infty,0)$上单调递减，在$(
      0,+\infty)$上单调递增.

      因$G(0)=-f(0)>0$，所以$G(x)>0$恒成立，故$F'(x)=\frac{G(x)}{x^{2}}>0$恒成立，

      即$F(x)$在$\left(-\infty,+\infty\right)$上单调递增.
  \end{enumerate}

  \section{综合题(每小题 5 分，共 10 分)}
  \begin{enumerate}
    \setcounter{enumi}{18}
    \item \textbf{Proof}. 由题可知$x_{1}=\sqrt{a}$，$x_{n+1}=\sqrt{a+x_{n}}$.

      显然$x_{1}=\sqrt{a}<\frac{1+\sqrt{1+4a}}{2}$.设$x_{k}\leq\frac{1+\sqrt{1+4a}}{2}$，则
      \[
        x_{k+1}=\sqrt{a+x_{k}}\leq\sqrt{a+\frac{1+\sqrt{1+4a}}{2}}=\frac{1+\sqrt{1+4a}}{2}
        .
      \]

      所以$\forall\,n\in\mathbf{N}^{*}$，$x_{n}\leq\frac{1+\sqrt{1+4a}}{2}$，即$\{
      x_{n}\}$有上界.

      又因为
      \[
        x_{n+1}-x_{n}=\sqrt{a+x_{n}}-x_{n}=\frac{a+x_{n}-x_{n}^{2}}{\sqrt{a+x_n}+x_{n}}
        =-\frac{(x_{n}-\frac{1+\sqrt{1+4a}}{2})(x_{n}-\frac{1-\sqrt{1+4a}}{2})}{\sqrt{a+x_n}+x_{n}}
        \geq0,
      \]

      所以$\{x_{n}\}$单调递增，故$\{x_{n}\}$收敛.

      设$\lim_{n\to\infty}x_{n}=A$，则$A=\sqrt{a+A}$，解得$A=\frac{1+\sqrt{1+4a}}{2}$.

    \item \textbf{Proof}. 假设$\forall\,x\in(0,2)$，$|f'(x)|<M$.

      由Lagrange中值定理，$\exists\,\eta,\delta\in\left(0,2\right)$，$|f(x)|=|f'(
      \eta)|x$，$|f(x)|=|f'(\delta)|(x-2)$.

      所以$|f(x)|<Mx$，$|f(x)|<M(2-x)$，相加得$|f(x)|<M$，这与$M=\max_{x\in(0,2)}
      \{|f(x)|\}$矛盾.

      所以$\exists\,\xi\in(0,2)$，使得$|f'(\xi)|\geq M$.
  \end{enumerate}
\end{document}



